\chapter{Horizonte demográfico}

\titular{El paquete proyecta las pensiones a 2031 y no declara un solo supuesto de población
para educación. \textbf{Discute el horizonte demográfico donde la transición empuja al alza y
guarda silencio donde empuja a la baja.} Con los denominadores puestos, el resultado es
nítido: las pensiones caen \textbf{4,8\,\%} por persona de 65 años y más, y la educación sube
\textbf{4,4\,\%} por persona en edad escolar.}

\quecambio{La población de 65 años y más creció 502\,822 personas en un año, un 4,3\,\%. La de
3 a 14 años cayó 248\,115 personas, un 1,0\,\%. \textbf{Los dos denominadores de la política
social se mueven en direcciones opuestas y el paquete solo reconoce uno.}}

\porqueimporta{Es el capítulo que la corrida anterior no pudo escribir. Se negó a calcular por
habitante porque los denominadores no estaban en el paquete, y \textbf{esa fue una decisión de
archivo, no de disciplina}: los denominadores existen, son públicos y están al alcance. La
diferencia entre tenerlos y no tenerlos es la diferencia entre decir que el gasto pensionario
crece y decir que la presión demográfica crece más rápido que él.}

% cuadro C14_2
\begin{table}[htbp]
\centering
\scriptsize
\caption{Los denominadores demográficos}
\label{cua:C14_2}
\begin{tabular}{p{7.0cm}rr}
\toprule
\textbf{Denominador} & \textbf{2025} & \textbf{2026} \\
\midrule
Población total & 133 367 428 & 134 407 258 \\
Población de 65 años y más & 11 698 499 & 12 201 321 \\
Población de 3 a 14 años & 25 669 536 & 25 421 421 \\
\bottomrule
\end{tabular}
\\[2pt]\begin{minipage}{0.92\textwidth}\scriptsize Fuente: CONAPO, proyecciones de población a mitad de año 1950--2070. Tier externa demográfica. Elaborado por el ITED.\end{minipage}
\end{table}

% cuadro C14_1
\begin{table}[htbp]
\centering
\scriptsize
\caption{Cifras por habitante y por denominador declarado}
\label{cua:C14_1}
\begin{tabular}{p{5.0cm}p{4.4cm}rrr}
\toprule
\textbf{Indicador} & \textbf{Denominador} & \textbf{Pesos 2025} & \textbf{Pesos 2026} & \textbf{Var. real \%} \\
\midrule
Salud, por habitante & poblacion total (CONAPO) & 6 662,2 & 7 248,9 & 3,8 \\
Educacion, por habitante & poblacion total (CONAPO) & 8 037,6 & 8 644,3 & 2,6 \\
Educacion, por persona de 3 a 14 anios & poblacion en edad escolar basica (CONAPO) & 41 759,6 & 45 704,1 & 4,4 \\
Pensiones y jubilaciones, por persona de 65 y mas & poblacion 65+ (CONAPO) & 139 989,3 & 139 670,1 & -4,8 \\
Ambiente y agua, por habitante & poblacion total (CONAPO) & 131,3 & 126,3 & -8,2 \\
IMSS-Bienestar, por persona del padron & padron IMSS-Bienestar 2T2025, 23 entidades & --- & 4 810,5 & --- \\
\bottomrule
\end{tabular}
\\[2pt]\begin{minipage}{0.92\textwidth}\scriptsize CUADRO SEPARADO de los presupuestales: ninguna cifra del paquete aparece aquí junto a un denominador externo sin marca. No se publica el gasto por afiliado a IMSS o ISSSTE ni por alumno matriculado: esos dos denominadores no se consiguieron. Tier externa demográfica. Fuente: CONAPO, proyecciones de población a mitad de año, y padrón de IMSS-Bienestar del segundo trimestre de 2025. Elaborado por el ITED.\end{minipage}
\end{table}


\section{Declaración de la capa}

\marginal{tier externa demográfica}
Todas las cifras de este capítulo combinan dos fuentes de naturaleza distinta y esa mezcla se
declara antes de la primera.

El numerador viene del paquete y tiene tier \texttt{oficial\_primaria}. El denominador viene de
las proyecciones de población de CONAPO y del padrón de IMSS-Bienestar, y tiene un tier
propio, \texttt{externa\_demografica}, que este documento estrena.

\textbf{Los cuadros presupuestales y los per cápita van separados}, y ninguna cifra del paquete
aparece junto a un denominador externo en la misma línea sin marca. Cada cifra per cápita dice
cuál es su denominador.

\section{La tríada, con sus denominadores}

Pensiones, salud y educación son las tres transferencias públicas de ciclo de vida. Comparten
tres cosas: son transferencias por edad, tienen un componente contributivo y uno no
contributivo, y \textbf{el paquete publica el flujo de las tres y el perfil de edad de
ninguna}.

Difieren en la dirección en que la demografía las empuja, y ahí está el hallazgo:

\begin{list}{}{\setlength{\leftmargin}{1.4em}\setlength{\itemsep}{3pt}}
\item[$\bullet$] \textbf{Pensiones.} Denominador natural: población de 65 años y más. Creció
\textbf{4,3\,\%} en un año. El gasto cayó 0,7\,\% real. \textbf{Por persona: $-4,8$\,\%.}
\textbf{El paquete proyecta este rubro a 2031}, de 4,4 a 4,7\,\% del PIB.
\item[$\bullet$] \textbf{Educación.} Denominador natural: matrícula, y en su defecto población
en edad escolar. Cayó \textbf{1,0\,\%} en un año. El gasto creció 3,4\,\% real. \textbf{Por
persona en edad escolar: $+4,4$\,\%.} \textbf{El paquete no declara ningún supuesto de
población para este rubro.}
\item[$\bullet$] \textbf{Salud.} Denominador mixto: la demanda crece por envejecimiento y el
perfil cubre todo el ciclo. El gasto creció 4,6\,\% real; \textbf{por habitante, 3,8\,\%.}
El paquete no proyecta el rubro.
\end{list}

\section{La asimetría, y qué habría significado corregirla}

\marginal{horizonte donde presiona al alza}
El marco de mediano plazo del CGPE proyecta pensiones y jubilaciones año por año hasta 2031.
No proyecta educación, ni salud, ni matrícula, ni cobertura.

\textbf{La asimetría no es casual y tiene una lógica presupuestal:} las pensiones son un
compromiso adquirido que hay que pagar, y proyectarlo es obligatorio para la sostenibilidad.
La educación es gasto discrecional anual y no obliga a proyectar nada.

\textbf{Pero la consecuencia es que el paquete no puede ver una economía que la demografía le
está dando.} Si la población en edad escolar cae 1,0\,\% al año y el paquete no lo declara, el
presupuesto educativo se discute como si el número de alumnos fuera constante. Con la capa
demográfica puesta, la pregunta cambia: \textbf{¿cuánto del aumento de 3,4\,\% real en
educación es mejor servicio y cuánto es simplemente menos gente que atender?}

Este documento puede acotar la respuesta y no cerrarla. \textbf{De los 4,4 puntos de aumento
por persona en edad escolar, aproximadamente uno viene del denominador} ---la caída de
248\,115 personas--- y el resto del numerador. Cerrarla exigiría la matrícula, que separa
demografía de cobertura, y \textbf{la matrícula no se obtuvo}.

\section{Dónde declara el paquete un supuesto demográfico}

Se buscó en las tres piezas del paquete. El resultado:

\begin{list}{}{\setlength{\leftmargin}{1.4em}\setlength{\itemsep}{2pt}}
\item[$\bullet$] \textbf{En pensiones, sí}, indirectamente: el marco de mediano plazo proyecta
el gasto pensionario a 2031 y esa proyección incorpora demografía, aunque no publique la
población que usa.
\item[$\bullet$] \textbf{En educación, no.} Ni matrícula, ni cobertura, ni población en edad
escolar, ni planteles, ni docentes.
\item[$\bullet$] \textbf{En salud, no.} Ni población derechohabiente, ni población sin
seguridad social, ni perfil de demanda por edad.
\end{list}

\textbf{El paquete asigna 3,8 billones de pesos a los tres rubros y no publica un solo
denominador para ninguno.}

\section{Implicaciones}

Tres, y la tercera es la que importa para el año que viene.

\textbf{Primera:} la presión pensionaria crece más rápido que el gasto que la financia. Cuatro
puntos y ocho décimas de caída real por persona de 65 años y más, en un ejercicio. Si eso se
repite, el problema deja de ser de sostenibilidad fiscal y pasa a ser de suficiencia de las
pensiones.

\textbf{Segunda:} la educación recibe un dividendo demográfico que nadie contabiliza. Cada año
hay menos personas en edad escolar, y cada año el gasto por persona mejora sin que se decida
mejorarlo. Es una oportunidad de política ---la misma plata alcanza para más por alumno--- y no
se puede aprovechar lo que no se mide.

\textbf{Tercera, sobre el propio documento:} este capítulo no se habría podido escribir hace un
año, y no por falta de datos sino por una regla mal aplicada. La regla de no publicar cifras
per cápita sin denominador declarado es correcta. \textbf{Convertirla en no calcular cifras per
cápita fue un error, y corregirlo costó dos descargas.}

\begin{ficha}{horizonte demográfico}
\fila{Perímetro}{Los tres agregados de la tríada, cada uno con el perímetro declarado en su
capítulo: pensiones por el Anexo 3, salud y educación por su función.}
\fila{Línea de comparación}{\textbf{Línea P} en los numeradores. Los denominadores son
proyecciones a mitad de año de cada ejercicio.}
\fila{Deflactor}{Del PIB, 4,8\,\%, CGPE 2026, aplicado al numerador antes de dividir.}
\fila{Año de los pesos}{Las cifras per cápita están en pesos corrientes de cada año; la
variación real usa el deflactor.}
\fila{PIB y añada}{38\,715,9 miles de millones, CGPE 2026, solo donde se cita razón a PIB.}
\fila{Denominador demográfico}{CONAPO, proyecciones de población a mitad de año 1950--2070,
por entidad, edad simple y sexo: población total, población de 65 años y más y población de 3 a
14 años. Padrón de personas sin seguridad social de IMSS-Bienestar, segundo trimestre de 2025,
veintitrés entidades.}
\fila{Fuentes}{CONAPO; IMSS-Bienestar; los cuadros presupuestales de los capítulos 5, 6 y 7;
CGPE 2026 Anexo III.2.}
\fila{Qué no puede afirmarse}{La matrícula por nivel educativo y por tanto la separación entre
efecto demográfico y efecto de cobertura. La población derechohabiente de IMSS y de ISSSTE. La
pensión media, que exige el padrón de pensionados. Y los perfiles de edad de los tres rubros,
que exigen cuentas de transferencias nacionales y quedan fuera del alcance de la lectura de un
paquete.}
\end{ficha}
\clearpage
